
% Compile this single file twice with pdflatex.
% All proof text, bibliography entries and the Python certificate are inline.
\documentclass[11pt,a4paper]{article}
\usepackage[margin=26mm]{geometry}
\usepackage[T1]{fontenc}
\usepackage{lmodern,amsmath,amssymb,amsthm,booktabs,array,longtable}
\usepackage{xcolor,listings}
\usepackage{tocloft}
\setlength{\cftbeforesecskip}{4pt}
\usepackage[hidelinks]{hyperref}
\usepackage{microtype}
\hypersetup{pdftitle={Exact Prime Covers and Arithmetic Obstructions},
 pdfauthor={Prepared for KEIO},pdfsubject={An integrated mathematical research manuscript}}
\newtheorem{theorem}{Theorem}[section]
\newtheorem{lemma}[theorem]{Lemma}
\newtheorem{proposition}[theorem]{Proposition}
\newtheorem{corollary}[theorem]{Corollary}
\theoremstyle{definition}\newtheorem{definition}[theorem]{Definition}
\theoremstyle{remark}\newtheorem{remark}[theorem]{Remark}
\newcommand{\Z}{\mathbb Z}\newcommand{\Q}{\mathbb Q}
\newcommand{\R}{\mathbb R}\newcommand{\C}{\mathbb C}
\newcommand{\N}{\mathbb N}
\newcommand{\Tr}{\operatorname{Tr}}\newcommand{\ord}{\operatorname{ord}}
\newcommand{\floor}[1]{\left\lfloor#1\right\rfloor}
\newcommand{\ceil}[1]{\left\lceil#1\right\rceil}
\newcommand{\ind}{\mathbf 1}
\setlength{\emergencystretch}{2em}
\setcounter{tocdepth}{1}
\title{\bfseries Exact Prime Covers and Arithmetic Obstructions\\[7pt]
\large Salem powers, repeated blocks, Euler integrals and Lambert functions}
\author{Prepared for KEIO\\[2pt]\small Integrated research manuscript}
\date{24 October 2026}
\begin{document}
\maketitle
\begin{abstract}
We assemble a source-audited collection of exact classifications and
arithmetic obstructions, retaining complete proofs and the original
target constants. For the largest root $\tau$ of
$X^4-X^3-X^2-X+1$, every $\floor{\tau^{hn}}$ has a proper divisor in
$\{2,3,5\}$ if and only if $260\mid h$; the natural density of terms
coprime to $30$ is calculated for every exponent. Further results
exclude macroscopic repetitions in three arithmetic words,
classify logarithmic contamination for rational floor kernels, and
show why the raw twin-prime Euler products have approximation exponent
zero relative to their reduced denominators. An asymmetric Euler
integral admits exact integer linear forms and short cross-index
logarithm elimination; equally short identically zero forms demonstrate
the unresolved nonvanishing requirement. Finally, a fixed-base Lambert
family representing the Erd\H{o}s--Borwein and reciprocal Fibonacci
constants has an exact all-jet orbit formula but admits no finite
rational power-Mahler system. The manuscript consolidates selected
HTML results with their corpus antecedents and the preceding Salem
classification. It claims neither historical priority nor a new
transcendence theorem for an unresolved named constant.
\end{abstract}
\tableofcontents
\newpage
\section{Purpose, selection and mathematical scope}\label{sec:introduction}
The starting materials are the supplied HTML research record
\cite{HTMLRecord}, the preceding fixed-exponent Salem note
\cite{SalemNote}, and the 49 TeX manuscripts in the repository snapshot
\cite{Corpus}. The selection is mathematical rather than chronological:
we retain results that give an exact classification, preserve an original
constant through a useful coordinate system, or isolate a demonstrable
failure of a proposed arithmetic argument. All selected statements are
proved here from their stated standard inputs.

Four different conclusions must be distinguished. A complete finite
classification can decide every exponent of a specified prime-cover
problem. A functional obstruction can exclude a specified system of
functional equations. An exact elimination identity can produce small
integer linear forms. None of the latter conclusions, by itself,
establishes transcendence of the value that motivated it. In particular,
a nonzero coefficient vector can produce the zero linear form.

The prime-cover classification in
Theorem~\ref{thm:main} and its general criterion in
Theorem~\ref{thm:criterion} constitute the main positive classification.
The other sections provide complementary information about repeated
arithmetic data, logarithmic contamination, denominator growth and
finite-dimensional functional methods. These are distinct mechanisms;
their integration is not a claim that one new principle solves all of
the underlying constant problems.

The comparison used the immutable commit
\begin{center}
\small\texttt{f70e5fedabf7b90313f703eb85f7dff72a055541}.
\end{center}
A result described as additional to that snapshot is not thereby
historically new. The source map in Appendix~\ref{app:sources} identifies
the retained HTML items, overlapping corpus statements and the
specific refinements included here. The repository describes its
contents as research drafts; source inclusion is not treated as
independent certification of a theorem. Likewise, the exact arithmetic
certificate in Appendix~\ref{app:certificate} verifies the finite
computations, while the analytic arguments in the text justify their
infinite-index consequences.
\section{The exact prime-cover classification}\label{sal:main-section}
A finite set of primes $S$ \emph{covers} a positive integer sequence if every term has a proper prime divisor in $S$. An \emph{eventual cover} allows finitely many initial exceptions. Covering is stronger than compositeness: a composite sequence need not admit any fixed finite cover.

The supplied manuscripts \cite{Corpus005,Corpus031} establish effective powers of Salem numbers with covered floors and classify the prime sets that work after choosing a suitable power. Here the exponent is prescribed, so the trace need not have its identity residue. The following specific classification is the principal positive classification of this paper.

\begin{theorem}\label{thm:main}
Let $\tau>1$ be the largest root of
\[
 P(X)=X^4-X^3-X^2-X+1.
\]
For every integer $h\ge1$, the following conditions are equivalent:
\begin{enumerate}
\item For every $n\ge1$, the integer $\floor{\tau^{hn}}$ has a proper prime divisor in $\{2,3,5\}$.
\item The same holds for all sufficiently large $n$.
\item $260\mid h$.
\end{enumerate}
The natural density
\[
 \delta(h)=\lim_{N\to\infty}\frac1N
 \#\{1\le n\le N:\gcd(30,\floor{\tau^{hn}})=1\}
\]
exists and is given by Table~\ref{tab:density}. In particular,
\[
 260\nmid h\quad\Longrightarrow\quad\delta(h)\ge\frac1{12},
\]
with equality precisely when $\gcd(h,260)=130$.
\end{theorem}

This makes $260$ the least exponent for the indicated three-prime cover. It does not make $260$ the least exponent for compositeness by an arbitrary mechanism, or for coverage by a larger finite prime set. In particular, positive density of terms coprime to $30$ is not a statement about prime terms.

The surrounding literature includes constructions of composite power floors and the explicit-example questions of Hahn--Ismailescu--Kim--Kim \cite{HIKK}. Those questions motivate the setting; this paper does not present their transcendental-number question as newly solved. 

\section{The conjugate correction and its distribution}
We use the convention that a Salem number has degree at least four. Let $\tau$ have degree $d=2s+2$, with conjugates
\[
 \tau,\tau^{-1},z_1,z_1^{-1},\ldots,z_s,z_s^{-1},
 \qquad |z_j|=1.
\]
Choose the $z_j$ one from each reciprocal pair. Define
\[
 T_m=\Tr_{\Q(\tau)/\Q}(\tau^m)\in\Z,
 \qquad E_m=\tau^{-m}+\sum_{j=1}^s(z_j^m+z_j^{-m}).
\]
The exact identity, including integral correction values, is
\begin{equation}\label{eq:floor}
 \floor{\tau^m}=T_m-\ceil{E_m}.
\end{equation}
We reproduce the needed analytic facts so that the finite calculation is not asked to stand in for the infinite argument.

\begin{lemma}[Strict correction window]\label{lem:window}
For every $m\ge1$,
\[
 -(d-2)<E_m<d-2,
 \qquad \ceil{E_m}\in Q_d:=\{-d+3,\ldots,d-2\}.
\]
\end{lemma}
\begin{proof}
None of the $z_j$ is a root of unity: otherwise the irreducible minimal polynomial of $\tau$ would be cyclotomic. Put $u=\tau^{-m}\in(0,1)$ and $x_j=|1-z_j^m|^2>0$. The nonzero algebraic-integer norm satisfies
\[
 1\le\left|N_{\Q(\tau)/\Q}(\tau^m-1)\right|
   =\frac{(1-u)^2}{u}\prod_{j=1}^sx_j.
\]
Thus $\prod_jx_j\ge u/(1-u)^2>u$. If $\sum_jx_j\le u$, the arithmetic--geometric mean inequality would give
\[
 \prod_jx_j\le(u/s)^s\le u,
\]
a contradiction. Since $E_m=u+2s-\sum_jx_j$, the upper bound follows. The lower bound follows from $E_m\ge u-2s>-2s$. Taking ceilings gives the stated finite integer set.
\end{proof}

\begin{lemma}[All corrections on every progression]\label{lem:distribution}
For each $q\in Q_d$, there is a real number $w_q>0$ such that, on every fixed positive-step arithmetic progression of $m$, the condition $\ceil{E_m}=q$ has relative natural density $w_q$. The weights are independent of the progression. In degree four,
\begin{equation}\label{eq:weights}
 (w_{-1},w_0,w_1,w_2)=\left(\frac13,\frac16,\frac16,\frac13\right).
\end{equation}
\end{lemma}
\begin{proof}
Write $z_j=e^{2\pi i\theta_j}$. First, $1,\theta_1,\ldots,\theta_s$ are linearly independent over $\Q$. More generally suppose $\prod_jz_j^{a_j}=\omega$, where the $a_j$ are integers and $\omega$ is a root of unity. For a fixed $j$, choose an automorphism of the splitting field sending $z_j$ to $\tau$. No other selected $z_k$ can map to $\tau$ by injectivity, or to $\tau^{-1}$ because that would force $z_k=z_j^{-1}$. Every other image therefore has modulus one. Taking absolute values after applying the automorphism yields $\tau^{a_j}=1$, so $a_j=0$. Repeating this proves the assertion.

For $m=an+b$, $a\ge1$, every nontrivial integral character of the vector $(m\theta_1,\ldots,m\theta_s)$ has irrational frequency. Its normalized exponential sum is a bounded geometric sum divided by its length and tends to zero. The Weyl criterion therefore gives equidistribution on the $s$-torus.

On that torus let
\[
 Y(t)=2\sum_{j=1}^s\cos(2\pi t_j),
 \qquad w_q=\operatorname{meas}\{t:q-1<Y(t)<q\}.
\]
The interval $(q-1,q)$ meets the interior of the range $[-2s,2s]$, and its inverse image contains a nonempty open set. Hence $w_q>0$. Each level set $Y=c$ has measure zero: after fixing $t_2,\ldots,t_s$, there are at most two possible values of $t_1$, so Fubini's theorem applies. Equidistribution now gives the frequencies. The vanishing term $\tau^{-m}$ does not change them, since neighborhoods of the two boundary level sets have measures tending to zero. In degree four the four arc lengths for $2\cos(2\pi t)$ give \eqref{eq:weights}.
\end{proof}

\section{A finite criterion for a fixed exponent}
Let $C$ be an integral companion matrix of the minimal polynomial of $\tau$. It is invertible modulo every positive integer, since a Salem number is a unit, and $T_m=\operatorname{tr}(C^m)$.

\begin{theorem}[Fixed-exponent criterion]\label{thm:criterion}
Let $S$ be a finite nonempty prime set, $M=\prod_{p\in S}p$, and $J\subset\Z$ finite and nonempty. Fix $h\ge1$ and put
\[
 L=\ord(C^h\bmod M),\qquad
 s_k=\operatorname{tr}(C^{hk})\bmod M\quad(0\le k<L).
\]
For all sufficiently large $n$, every $\floor{\tau^{hn}}+j$, $j\in J$, has a proper prime divisor in $S$ if and only if
\begin{equation}\label{eq:criterion}
 \gcd(M,s_k-q+j)>1
 \quad(0\le k<L,\ q\in Q_d,\ j\in J).
\end{equation}
For each fixed $j$, the natural density of indices with
$\gcd(M,\floor{\tau^{hn}}+j)=1$ is
\begin{equation}\label{eq:densitygeneral}
 D_{h,j}=\frac1L\sum_{k=0}^{L-1}\sum_{q\in Q_d}
 w_q\ind_{\gcd(M,s_k-q+j)=1}.
\end{equation}
\end{theorem}
\begin{proof}
If \eqref{eq:criterion} holds, \eqref{eq:floor} and Lemma~\ref{lem:window} give a prime divisor in $S$ for every term. For sufficiently large $n$, every term is larger than $\max S$, so the divisor is proper.

If the condition fails at $(k,q,j)$, restrict to positive indices $n\equiv k\pmod L$. The trace is $s_k$ modulo $M$. Lemma~\ref{lem:distribution} supplies a positive relative-density subset on which $\ceil{E_{hn}}=q$. On that subset \eqref{eq:floor} makes the shifted floor coprime to $M$, contradicting eventual coverage. Finally, split all indices into their $L$ residue progressions and apply the same frequency statement to each; the resulting finite sum is \eqref{eq:densitygeneral}.
\end{proof}

\begin{corollary}[Finite search over exponents]\label{cor:search}
Put $o=\ord(C\bmod M)$. The criterion and densities of Theorem~\ref{thm:criterion} depend on $h$ only through $g=\gcd(h,o)$. If any exponent gives an eventual cover, the least such exponent is the smallest divisor $g$ of $o$ whose trace orbit passes \eqref{eq:criterion}.
\end{corollary}
\begin{proof}
The cyclic subgroups generated by $C^h$ and $C^g$ are equal. Their elements, and hence their traces with multiplicity, differ only by a permutation. Moreover $g\le h$ and $g$ itself is an admissible exponent whenever $h$ is.
\end{proof}

A smaller scalar period of $(T_m\bmod M)$ may replace the matrix order in a particular example. The quartic below has scalar period $260$, although its companion matrix has order $780$ modulo $30$.

\section{The quartic: exact arithmetic and proof of the main theorem}
The polynomial $P$ in Theorem~\ref{thm:main} is irreducible modulo two. It has no root in $\mathbb F_2$, and it is not $(X^2+X+1)^2$, the only possible product of irreducible quadratic polynomials in degree four over that field. Thus $P$ is irreducible over $\Q$.

Dividing by $X^2$ and putting $y=X+X^{-1}$ gives $y^2-y-3$. Its roots are $(1\pm\sqrt{13})/2$, one above $2$ and one in $(-2,2)$. Consequently $\tau$ is a quartic Salem number. Comparing $\tau+\tau^{-1}=(1+\sqrt{13})/2$ with $3/2+2/3$ also gives $\tau>3/2$.

The traces satisfy
\begin{equation}\label{eq:recurrence}
 \begin{gathered}
 T_0=4,\quad T_1=1,\quad T_2=3,\quad T_3=7,\\
 T_{m+4}=T_{m+3}+T_{m+2}+T_{m+1}-T_m.
 \end{gathered}
\end{equation}
The integer calculations used below are reproduced by the standard-library Python certificate embedded in Appendix~\ref{app:certificate}. They can also be obtained directly from \eqref{eq:recurrence}. No floating-point approximation is used.

The four-term state modulo $30$ returns to its initial state at $260$. It does not return at $260/p$ for $p=2,5,13$. The recurrence transition is invertible modulo $30$, because its trailing coefficient is $-1$. These tests prove that the scalar period is exactly $260$: the period divides $260$, and any proper divisor would divide $260/p$ for at least one prime divisor $p$ of $260$. An independent matrix calculation verifies order $780$ for
\[
 C=\begin{pmatrix}
 0&0&0&-1\\1&0&0&1\\0&1&0&1\\0&0&1&1
 \end{pmatrix}\pmod{30}.
\]
Specifically, $C^{780}=I$ and $C^{780/p}\ne I$ for $p=2,3,5,13$.

For $Q_4=\{-1,0,1,2\}$, define a trace residue $t$ to be safe when
\[
 \gcd(30,t-q)>1\quad\hbox{for every }q\in Q_4.
\]
The safe residues modulo $30$ are exactly $4,5,26,27$. Every proper divisor $g$ of $260$ has an unsafe trace $T_g$, as shown in Table~\ref{tab:unsafe}.

\begin{table}[ht]
\centering\small
\begin{tabular}{@{}r|rrrrrrrrrrr@{}}
\toprule
$g$&1&2&4&5&10&13&20&26&52&65&130\\
\midrule
$T_g\bmod30$&1&3&7&16&18&1&22&9&19&16&24\\
\bottomrule
\end{tabular}
\caption{All proper divisors of the trace period yield unsafe traces.}
\label{tab:unsafe}
\end{table}

\begin{proof}[Proof of Theorem~\ref{thm:main}, equivalences]
If $260\nmid h$, put $g=\gcd(h,260)<260$. There is an integer $n_0$ with $hn_0\equiv g\pmod{260}$, because $h/g$ is invertible modulo $260/g$. Along the corresponding progression of positive $n$, the trace is $T_g$ modulo $30$. Table~\ref{tab:unsafe} supplies an unsafe trace; Lemma~\ref{lem:distribution} supplies a positive relative density of a correction $q$ making the floor coprime to $30$. Eventual coverage therefore fails.

If $260\mid h$, every trace $T_{hn}$ equals $T_0=4$ modulo $30$. The correction window gives floor residues $2,3,4,5$, all divisible by one of $2,3,5$. Since $h\ge260$ and $\tau>3/2$, every floor is greater than $5$, making the divisor proper for every $n\ge1$. This proves all three equivalences.
\end{proof}

\section{The exact density classification}
Let $g=\gcd(h,260)$ and $\ell=260/g$. Multiplication by $h/g$ permutes the residue classes modulo $\ell$. Combining the scalar trace period with \eqref{eq:weights} and \eqref{eq:densitygeneral} gives
\begin{equation}\label{eq:densityquartic}
 \delta(h)=\frac1\ell\sum_{a=0}^{\ell-1}\sum_{q=-1}^{2}
 w_q\ind_{\gcd(30,T_{ga}-q)=1}.
\end{equation}
This is a finite sum of rational numbers, and yields Table~\ref{tab:density}.

\begin{table}[ht]
\centering
\begin{tabular}{@{}rr@{\qquad\qquad}rr@{}}
\toprule
$\gcd(h,260)$&$\delta(h)$&$\gcd(h,260)$&$\delta(h)$\\
\midrule
1&$541/1560$&20&$8/39$\\
2&$241/780$&26&$7/20$\\
4&$14/39$&52&$2/5$\\
5&$97/312$&65&$5/24$\\
10&$15/52$&130&$1/12$\\
13&$49/120$&260&$0$\\
\bottomrule
\end{tabular}
\caption{Natural density of power floors coprime to $30$, for every exponent.}
\label{tab:density}
\end{table}

All nonzero entries are at least $1/12$, and only the row $g=130$ attains that bound. This completes the density assertion of Theorem~\ref{thm:main}. For instance $h=130$ fails the three-prime cover on exactly $1/12$ of indices, whereas $h=260$ has the cover on every index.

The certificate checks \eqref{eq:densityquartic} by exact rational arithmetic. It independently computes the same densities using the matrix orbit of length dividing $780$, and verifies agreement for every exponent residue class modulo $780$. The analytic proof is what converts these finite orbit sums into natural densities; the certificate alone does not establish equidistribution.


% HTML-specific arithmetic supplement.  Labels use the arith: prefix.
% Required bibliography keys: Corpus001, Corpus003, ArithDavenport, ArithHW.
% Uses only the theorem, lemma, corollary, remark environments in the main file.

\section{Macroscopic repetitions in three arithmetic words}
\label{arith:sec:repetition}

This section retains the coefficient sequences of the original constants.
It adds the repetition statements of HTML T1--T3 to the prime-parity coding
already present in the corpus \cite{Corpus003}.  The tools are classical:
the prime number theorem, Siegel--Walfisz, and the divisor bound.  No claim
of a new transcendence principle or historical priority is made.

For a sequence $w=(w_n)_{n\ge1}$, an \emph{equal-block triple} is a triple
of integers $(r,s,L)$ with $0\le r<s$ and $L\ge1$ such that
\[
 w_{r+i}=w_{s+i}\qquad(1\le i\le L).
\]
We say that $w$ has \emph{no macroscopic repetitions} if, for every
$\varepsilon>0$, there is $s_0(\varepsilon)$ such that no equal-block
triple has $s\ge s_0(\varepsilon)$ and $L\ge\varepsilon s$.
This definition concerns exact equality of the indicated coefficient
blocks.  It does not assert that the sequence has no short repetitions.

Let $p_j$ be the $j$th prime, let $\pi(x)$ count primes at most $x$, and
let $d(n)$ count the positive divisors of $n$.

\begin{theorem}[Three repetition obstructions]\label{arith:thm:words}
Each of the following sequences has no macroscopic repetitions:
\begin{enumerate}
\item $b_n^{(m)}=\pi(n)\bmod m$, with residues in $\{0,\ldots,m-1\}$,
for every fixed integer $m\ge2$;
\item the prime-gap sequence $g_n=p_{n+1}-p_n$;
\item the divisor-count sequence $d(n)$.
\end{enumerate}
\end{theorem}

We first isolate the exact analytic input in the first two assertions.
The classical Siegel--Walfisz theorem states that, for every fixed
$A,B>0$, uniformly for $q\le(\log X)^A$ and $(a,q)=1$,
\begin{equation}\label{arith:eq:SW}
 \pi(X;q,a)=\frac{\operatorname{Li}(X)}{\varphi(q)}
       +O_{A,B}\!\left(\frac{X}{(\log X)^B}\right).
\end{equation}
Here $\pi(X;q,a)$ counts primes in the residue class $a$ modulo $q$.
The constants need not be effective for the arguments below.
We also use $\vartheta(Y)=\sum_{p\le Y}\log p\sim Y$ and
$p_n\sim n\log n$.  These are standard prime-number theorems; see
\cite{ArithDavenport}.

\begin{lemma}[A prime in the necessary reduced class]\label{arith:lem:prime}
Fix $0<a<b$ and $C>0$.  For every sufficiently large $X$, every prime
$\ell\le C\log X$, and every integer $D$ with $\ell\nmid D$, there is a
prime in $[aX,bX]$ congruent to $D$ modulo $\ell$.
Moreover, if $1\le D\le2X$, then some prime $\ell\ll\log X$ does not
divide $D$; the implied constant is absolute.
\end{lemma}
\begin{proof}
Apply \eqref{arith:eq:SW}, with $A=2$ and $B=5$, at $aX$ and $bX$.
For large $X$ its modulus range includes every $\ell\le C\log X$.
The difference of the main terms is
\[
 \frac{\operatorname{Li}(bX)-\operatorname{Li}(aX)}{\ell-1}
       \gg_{a,b,C}\frac{X}{(\log X)^2},
\]
whereas the total error is $O_{a,b}(X/(\log X)^5)$.  The count is
therefore positive.  For the last assertion, if every prime at most
$4\log X$ divided $D$, then
\[
 D\ge\prod_{p\le4\log X}p
      =\exp(\vartheta(4\log X))>X^2>2X
\]
for large $X$, a contradiction.
\end{proof}

\begin{proof}[Proof of Theorem~\ref{arith:thm:words}, part (1)]
Suppose that $b^{(m)}$ has equal blocks with $L\ge\varepsilon s$ and
arbitrarily large $s$.  Put $D=s-r$.  Adjacent differences recover the
prime indicator:
\[
 b_n^{(m)}-b_{n-1}^{(m)}\equiv\mathbf1_{\{n\text{ prime}\}}\pmod m.
\]
Since an indicator is either $0$ or $1$ and $m\ge2$, equality of adjacent
residue differences implies equality of the indicators.  Thus
\begin{equation}\label{arith:eq:translated-indicators}
 \mathbf1_{\{r+i\text{ prime}\}}
   =\mathbf1_{\{s+i\text{ prime}\}}\qquad(2\le i\le L).
\end{equation}
Choose $c=\min(\varepsilon/4,1/4)>0$.  The interval
$[cs,2cs]$ lies in $[2,L]$ for large $s$.  As $1\le D\le s$, choose by
Lemma~\ref{arith:lem:prime} a prime $\ell\ll\log s$ not dividing $D$.
The same lemma gives a prime
\[
 y\in[(1+c)s,(1+2c)s],\qquad y\equiv D\pmod\ell.
\]
Write $i=y-s$.  Then $2\le i\le L$, while
$r+i=y-D\ge cs>\ell$ and $r+i\equiv0\pmod\ell$.
Consequently $r+i$ is composite and $s+i=y$ is prime, contradicting
\eqref{arith:eq:translated-indicators}.
\end{proof}

\begin{proof}[Proof of Theorem~\ref{arith:thm:words}, part (2)]
Suppose that $g_{r+i}=g_{s+i}$ for $1\le i\le L$, where
$L\ge\varepsilon s$ and $s$ is arbitrarily large.  Shortening both
blocks if necessary, take $L=\lceil\varepsilon s\rceil$.
Put $X=p_{s+1}$ and $D=p_{s+1}-p_{r+1}$, so $0<D<X$.
Summing the equal gaps gives
\begin{equation}\label{arith:eq:translated-primes}
 p_{r+j}+D=p_{s+j}\qquad(1\le j\le L+1).
\end{equation}
The two lists contain all primes in their respective endpoint intervals.
Translation by $D$ therefore preserves primality for every integer in
$[p_{r+1},p_{r+L+1}]$.

The prime number theorem gives
$p_{s+L+1}/X\longrightarrow1+\varepsilon$.
With $c=\min(\varepsilon/8,1/8)$, the interval
$[(1+c)X,(1+2c)X]$ lies strictly inside
$[p_{s+1},p_{s+L+1}]$ for large $s$.
Choose a prime $\ell\ll\log X$ not dividing $D$, and then a prime
$y$ in that interior interval with $y\equiv D\pmod\ell$, using
Lemma~\ref{arith:lem:prime}.  Its translate satisfies
\[
 y-D\ge p_{r+1}+cX>\ell,\qquad y-D\equiv0\pmod\ell.
\]
It is composite, contradicting the primality preservation just proved.
\end{proof}

\begin{proof}[Proof of Theorem~\ref{arith:thm:words}, part (3)]
Suppose that divisor counts agree on two blocks as above, and shorten
them to length $L=\lceil\varepsilon s\rceil$.  Put $h=s-r$, so
$1\le h\le s$.  The elementary identity
\[
 d(n)\equiv1\pmod2\quad\Longleftrightarrow\quad n\text{ is a square}
\]
shows that every square $k^2$ in $[r+1,r+L]$ translates to a square
$k^2+h=j^2$.
The number of squares in this interval is at least
\[
 \sqrt{r+L}-\sqrt r-1
 =\frac{L}{\sqrt{r+L}+\sqrt r}-1
 \gg_{\varepsilon}\sqrt s.
\]
On the other hand, each such pair satisfies
$(j-k)(j+k)=h$.  There are at most $d(h)$ possibilities, since the
positive factor $j-k$ determines $j+k$ and hence $k$.
The classical divisor bound $d(n)=O_{\eta}(n^{\eta})$ for every
$\eta>0$ \cite{ArithHW}, with $\eta=1/4$, gives
$d(h)=O(s^{1/4})$.  This contradicts the displayed lower bound.
\end{proof}

\begin{remark}[What the coefficient obstruction does and does not say]
The natural generating series retain the original constants:
\[
 2\sum_{k\ge1}(-1)^{k+1}2^{-p_k}
       =\sum_{n\ge1}b_n^{(2)}2^{-n},\qquad
 \sum_{n\ge1}\frac1{2^n-1}=\sum_{n\ge1}d(n)2^{-n}.
\]
The first follows by summation of adjacent differences, and the second
by absolutely convergent rearrangement of a geometric series.
The divisor coefficients are not binary digits: carrying can change
block structure.  Theorem~\ref{arith:thm:words} excludes sufficient
conditions requiring the specified long exact coefficient repetitions.
It proves neither transcendence nor its impossibility by other methods.
No assertion about a refined Diophantine exponent is needed here.
\end{remark}

\section{A classification of logarithmic contamination}
\label{arith:sec:poles}

For a rational function $g$ satisfying $g(u)=O(u^{-3})$ at infinity,
with poles only at negative integers, put
\[
 \mathcal J(g)=\int_0^{\infty}(\lfloor u\rfloor+1)g(u)\,du.
\]
The corpus already gives the complete pole-order evaluation of this
transform \cite{Corpus001}.  HTML T9 adds a classification for simple
poles and for two pole locations.  We reproduce the needed transform
calculation to make that classification independent of the longer source.
Write $H_0=0$ and $H_n=\sum_{j=1}^n1/j$.

\begin{lemma}[Simple and double poles]\label{arith:lem:transform}
Suppose
\[
 g(u)=\sum_{p\ge1}\left(\frac{\alpha_p}{u+p}
                         +\frac{\beta_p}{(u+p)^2}\right),
 \qquad \alpha_p,\beta_p\in\mathbb Q,
\]
with finitely many nonzero coefficients and $g(u)=O(u^{-3})$.
Then
\begin{align}
 \sum_p\alpha_p&=0,&\sum_p\beta_p&=\sum_p p\alpha_p,
       \label{arith:eq:cancellation}\\
 \mathcal J(g)&=\mathcal L(g)+
                   \sum_p\beta_p(\gamma-H_{p-1}),&
 \mathcal L(g)&=\sum_p\alpha_p\log\Gamma(p).
       \label{arith:eq:transducer}
\end{align}
Here $\gamma$ is Euler's constant.  In this specified decomposition the
coefficient of $\gamma$ is $B(g)=\sum_p\beta_p$.
\end{lemma}
\begin{proof}
The coefficients of $u^{-1}$ and $u^{-2}$ at infinity give
\eqref{arith:eq:cancellation}.  For a positive integer $N$, elementary
summation over the intervals $[j,j+1]$ gives
\begin{align*}
 A_p(N)&:=\int_0^N\frac{\lfloor u\rfloor+1}{u+p}\,du
       =N\log(N+p)-\log\Gamma(N+p)+\log\Gamma(p),\\
 B_p(N)&:=\int_0^N\frac{\lfloor u\rfloor+1}{(u+p)^2}\,du
       =H_{N+p-1}-H_{p-1}-\frac{N}{N+p}.
\end{align*}
Stirling's formula and $H_N=\log N+\gamma+o(1)$ yield
\begin{align*}
 A_p(N)&=N+(1/2-p)\log N+p
             -\tfrac12\log(2\pi)+\log\Gamma(p)+o(1),\\
 B_p(N)&=\log N+\gamma-H_{p-1}-1+o(1).
\end{align*}
Multiply by the partial-fraction coefficients and add.
The divergent terms cancel by \eqref{arith:eq:cancellation}, and
$\sum p\alpha_p=\sum\beta_p$ cancels the remaining constant terms.
Since the original weighted integral is absolutely convergent,
taking $N\to\infty$ proves \eqref{arith:eq:transducer}.
\end{proof}

\begin{theorem}[Simple poles and the complete two-location case]
\label{arith:thm:pole-classification}
The following statements hold.
\begin{enumerate}
\item Let $d\ge3$, $1\le p_1<\cdots<p_d$ be integers, and $C>0$ be
rational.  For
\[
 g(u)=\frac{C}{\prod_{i=1}^d(u+p_i)},
\]
one has $B(g)=0$ and
\[
 \mathcal L(g)>\frac{C}{(d-1)p_d^{\,d-1}}>0.
\]
\item Suppose that the poles of $g$ occur at exactly the two locations
$-a,-b$, where $1\le a<b$ are integers, each with multiplicity at most
two, and that $g(u)=O(u^{-3})$.  Then
\begin{equation}\label{arith:eq:two-poles}
 \mathcal L(g)=
    \frac{\log\Gamma(b)-\log\Gamma(a)}{b-a}\,B(g).
\end{equation}
Consequently, when $B(g)\ne0$, its logarithmic part vanishes if and only
if $(a,b)=(1,2)$.
\end{enumerate}
\end{theorem}
\begin{proof}
For (1), all $\beta_p$ vanish and
$\alpha_{p_i}=C/\prod_{j\ne i}(p_j-p_i)$.
Let $f(x)=\log\Gamma(x)$ and write $[p_1,\ldots,p_d]f$ for its divided
difference.  The mean-value theorem for divided differences gives
some $\xi\in(p_1,p_d)$ such that
\[
 \mathcal L(g)=C(-1)^{d-1}[p_1,\ldots,p_d]f
       =\frac{C(-1)^{d-1}}{(d-1)!}f^{(d-1)}(\xi).
\]
For every integer $m\ge2$ the convergent polygamma series is
\[
 f^{(m)}(x)=(-1)^m(m-1)!\sum_{n\ge0}(n+x)^{-m},\qquad x>0.
\]
Thus
\[
 \mathcal L(g)=\frac{C}{d-1}\sum_{n\ge0}(n+\xi)^{-(d-1)}
       >\frac{C}{(d-1)p_d^{d-1}},
\]
which proves the claim.  These are classical identities; they also follow
by differentiating the convergent product for $\Gamma$.

For (2), write the coefficients at the two locations as
$\alpha_a,\alpha_b,\beta_a,\beta_b$.  Equations
\eqref{arith:eq:cancellation} give
\[
 \alpha_b=-\alpha_a,\qquad
 B(g)=a\alpha_a+b\alpha_b=(b-a)\alpha_b.
\]
Substitution into the definition of $\mathcal L(g)$ proves
\eqref{arith:eq:two-poles}.  At positive integer arguments,
$\Gamma(1)=\Gamma(2)=1$ and $\Gamma(n)=(n-1)!$ is strictly increasing
for $n\ge2$.  The equality $\Gamma(a)=\Gamma(b)$ with $a<b$
therefore holds precisely for $(a,b)=(1,2)$.
\end{proof}

\begin{remark}[Scope of the classification]\label{arith:rem:pole-scope}
The words ``coefficient'' and ``logarithmic part'' refer to the explicit
partial-fraction decomposition, not to a claimed linear independence
of $1$, $\gamma$, and logarithms.  The exceptional pair gives, for example,
\[
 \mathcal J\!\left(\frac1{(u+1)^2(u+2)}\right)=\gamma,
 \qquad
 \mathcal J\!\left(\frac1{(u+1)(u+2)^2}\right)=1-\gamma.
\]
These are exact representations.  They do not give a sequence of
increasingly accurate rational approximations.  Within the class of
multiplicity at most two, a logarithm-free kernel with nonzero
$B(g)$ and pole locations other than the fixed pair $\{1,2\}$ needs
at least three locations and at least one double pole.  Higher pole
orders and determinant constructions lie outside this classification.
\end{remark}

\section{The raw twin-prime Euler-product approximation}
\label{arith:sec:twin}

HTML T13 concerns the reduced denominator of the original finite Euler
product.  Its conclusion is an obstruction to that specific approximation
sequence, independent of the twin-prime conjecture.
Define
\[
 C_2=\prod_{p>2}\left(1-\frac1{(p-1)^2}\right)>0,
 \qquad C_2(x)=\prod_{2<p\le x}\left(1-\frac1{(p-1)^2}\right).
\]
Positivity and convergence follow from
$\sum_{p>2}(p-1)^{-2}<\infty$ and positivity of every factor.
For $x\ge3$ write $C_2(x)=A_x/Q_x$ in lowest terms, with $Q_x>0$.

\begin{theorem}[Vanishing denominator-normalized exponent]
\label{arith:thm:twin}
As $x\to\infty$,
\begin{gather}
 Q_x\ge4^{\pi(x)-1},\qquad
 C_2(x)-C_2\sim\frac{C_2}{x\log x},
       \label{arith:eq:twin-bounds}\\
 \frac{-\log|C_2(x)-C_2|}{\log Q_x}\longrightarrow0,
 \qquad Q_x^{\eta}|C_2(x)-C_2|\longrightarrow\infty
       \quad\text{for every }\eta>0.
       \label{arith:eq:twin-exponent}
\end{gather}
\end{theorem}
\begin{proof}
Every numerator $p(p-2)$ in
$C_2(x)=\prod_{2<p\le x}p(p-2)/(p-1)^2$ is odd.
Reduction of the product therefore removes none of the powers of two
in its denominator.  With $v_2$ the usual valuation,
\[
 v_2(Q_x)=2\sum_{2<p\le x}v_2(p-1)\ge2(\pi(x)-1),
\]
which proves the denominator bound.

For completeness, partial summation and the prime number theorem give
\begin{align*}
 \sum_{p>x}\frac1{p^2}
  &=-\frac{\pi(x)}{x^2}+2\int_x^{\infty}\frac{\pi(t)}{t^3}\,dt\\
  &\sim-\frac1{x\log x}
                  +2\int_x^{\infty}\frac{dt}{t^2\log t}
   \sim\frac1{x\log x}.
\end{align*}
Indeed the relative prime-number-theorem error is uniformly small for
$t\ge x$ once $x$ is large, and integration by parts gives
$\int_x^{\infty}dt/(t^2\log t)\sim1/(x\log x)$.
Replacing $p^{-2}$ by $(p-1)^{-2}$ changes the sum by
$O(\sum_{n>x}n^{-3})=O(x^{-2})$.  Therefore
\[
 R_x:=-\log\prod_{p>x}\left(1-\frac1{(p-1)^2}\right)
 =\sum_{p>x}\frac1{(p-1)^2}
       +O\!\left(\sum_{p>x}p^{-4}\right)
 \sim\frac1{x\log x}.
\]
Since $C_2(x)=C_2e^{R_x}$, the error estimate in
\eqref{arith:eq:twin-bounds} follows.

Finally, $\log Q_x\ge(\pi(x)-1)\log4\gg x/\log x$, whereas
\[
 -\log|C_2(x)-C_2|=\log x+\log\log x+O(1).
\]
Their ratio tends to zero.  For fixed $\eta>0$, the same bounds give
\[
 \log\bigl(Q_x^{\eta}|C_2(x)-C_2|\bigr)
 \ge \eta(\pi(x)-1)\log4-\log x-\log\log x+O(1)
 \longrightarrow\infty,
\]
proving the final assertion.
\end{proof}

\begin{remark}[The route excluded]
If $C_2=u/v$ were rational, rational separation would give
$|C_2(x)-C_2|\ge1/(vQ_x)$; the difference is nonzero since the omitted
Euler factors are strictly below one.  Theorem~\ref{arith:thm:twin}
shows that the actual error is much larger than this lower bound.
The same observation applies to any fixed power of the denominator.
It does not decide the arithmetic nature of $C_2$ and does not exclude
weighted products or combinations of several truncations.
\end{remark}

\section{Cross-index elimination for an asymmetric Euler integral}
\label{comp:section}

This section preserves a useful synthesis from the HTML research record.
The exact integral family and its logarithmic obstruction are already in
manuscript~[001] \cite{Corpus001}, while the moving-window counting mechanism is already in manuscript~[021] \cite{Corpus021}.  Combining them gives the
following precise statement, with the nonvanishing issue left explicit.
It is not presented as a new proof mechanism or as an irrationality
result for Euler's constant.

For $n\geq1$, define
\begin{equation}\label{comp:integral}
 J_n=\int_0^1\!\int_0^1
 \frac{[x(1-x)y(1-y)]^n(1-x)}
 {(1-xy)(-\log(xy))}\,dx\,dy,
 \qquad d_n=\operatorname{lcm}(1,\ldots,2n+2).
\end{equation}
Write $H_k=\sum_{j=1}^k1/j$, $H_0=0$, and let $\gamma$ denote
Euler's constant.  All logarithms below are natural.

\begin{theorem}[Exact form and cross-index elimination]
\label{comp:main}
There are integers $a_n>0$, $b_n$, and $c_{n,\ell}\geq0$, indexed by
primes $\ell\leq2n+1$, such that
\begin{equation}\label{comp:integer-form}
 d_nJ_n=a_n\gamma+b_n+
 \sum_{\substack{\ell\leq2n+1\\\ell\ {\rm prime}}}
 c_{n,\ell}\log\ell,
 \qquad a_n=d_n\binom{2n+1}{n}.
\end{equation}
The logarithmic sum is strictly positive.  With $c=\log16-2>0$,
\begin{align}
 0<d_nJ_n&\leq\exp(-cn+o(n)),\label{comp:smallness}\\
 \log\max\bigl(1,a_n,|b_n|,\max_\ell c_{n,\ell}\bigr)
 &\leq(2+\log4)n+o(n).\label{comp:height}
\end{align}
For every sufficiently large $N$, there is a nonzero integer vector
$(u_N,\ldots,u_{2N})$ with
\begin{equation}\label{comp:vector-height}
 \log\max\bigl(1,\max_{N\leq n\leq2N}|u_n|\bigr)
 =O(N/\log N)
\end{equation}
which eliminates every logarithm in the moving window:
\begin{equation}\label{comp:compressed}
 \sum_{n=N}^{2N}u_nd_nJ_n=Q_N\gamma+P_N,
 \quad P_N,Q_N\in\mathbb Z,
 \quad |Q_N\gamma+P_N|\leq\exp(-cN+o(N)).
\end{equation}
Every nonzero eliminating vector has both positive and negative entries.
Moreover, there also exist nonzero vectors satisfying
\eqref{comp:vector-height} for which $P_N=Q_N=0$ exactly.
Consequently, the assertion that the vector is nonzero cannot be
substituted for the assertion that its resulting linear form is nonzero.
\end{theorem}

\subsection{Exact coefficients from a half-line kernel}

The identities
\[
 \frac1{-\log(xy)}=\int_0^\infty(xy)^t\,dt,
 \qquad \frac1{1-xy}=\sum_{j\geq0}(xy)^j
\]
have nonnegative integrands.  Tonelli's theorem, followed by the change
of variable $u=t+j$, therefore gives
\begin{align}
 J_n&=\int_0^\infty(\lfloor u\rfloor+1)g_n(u)\,du,
 \label{comp:half-line}\\
 g_n(u)&=B(u+n+1,n+2)B(u+n+1,n+1)
 =\frac{n!(n+1)!}{(u+M)\prod_{p=L}^{M-1}(u+p)^2},
 \label{comp:kernel}
\end{align}
where $L=n+1$, $M=2n+2$, and
$B(s,t)=\int_0^1x^{s-1}(1-x)^{t-1}\,dx$.
The factor $\lfloor u\rfloor+1$ counts the nonnegative integers $j\leq u$;
its values at integer endpoints do not affect the integral.

The partial-fraction expansion is
\begin{equation}\label{comp:partial-fractions}
 g_n(u)=\sum_{p=L}^{M}\frac{\alpha_p}{u+p}
       +\sum_{p=L}^{M-1}\frac{\beta_p}{(u+p)^2},
\end{equation}
where
\begin{align}
 \beta_{L+k}&=\binom nk\binom{n+1}{k},\qquad 0\leq k\leq n,
 \label{comp:beta}\\
 \alpha_{L+k}&=\beta_{L+k}
 \left(2H_k-2H_{n-k}-\frac1{n+1-k}\right),
 \qquad \alpha_M=\frac1{n+1}.
 \label{comp:alpha}
\end{align}
Indeed, evaluating $(u+L+k)^2g_n(u)$ at $u=-L-k$ gives
\[
 \beta_{L+k}=
 \frac{n!(n+1)!}{(n+1-k)(k!)^2((n-k)!)^2},
\]
and logarithmic differentiation gives the corresponding $\alpha_{L+k}$.
The remaining simple residue gives $\alpha_M$.  Comparing the first two
terms at infinity in \eqref{comp:partial-fractions} gives
\begin{equation}\label{comp:cancellations}
 \sum_p\alpha_p=0,
 \qquad\sum_p p\alpha_p=\sum_p\beta_p.
\end{equation}

Applying the common-cutoff transform in
Lemma~\ref{arith:lem:transform} to \eqref{comp:partial-fractions} gives
\begin{equation}\label{comp:exact-unscaled}
 J_n=\binom{2n+1}{n}\gamma
      -\sum_{p=L}^{M-1}\beta_pH_{p-1}
      +\sum_{p=L}^{M}\alpha_p\log((p-1)!).
\end{equation}
Here Vandermonde's identity gives
$\sum_{k=0}^n\binom nk\binom{n+1}{k}=\binom{2n+1}{n}$.
In particular, the coefficient of $\gamma$ is exact, not asymptotic.

Let $v_\ell(m)$ be the exponent of a prime $\ell$ in the positive
integer $m$.  Explicit coefficients for \eqref{comp:integer-form} are
\begin{align}
 b_n&=-d_n\sum_{p=L}^{M-1}\beta_pH_{p-1},
 \label{comp:b}\\
 c_{n,\ell}
 &=d_n\sum_{p=L}^{M}\alpha_pv_\ell((p-1)!)
 =d_n\sum_{j=L}^{M-1}T_jv_\ell(j),
 \qquad T_j=\sum_{p=j+1}^{M}\alpha_p.
 \label{comp:c}
\end{align}
The last equality follows by expanding the factorial logarithms and
using $\sum_p\alpha_p=0$ for the terms $j<L$.
Every denominator in $\alpha_p$ divides
$\operatorname{lcm}(1,\ldots,n+1)$, and every denominator in
$H_{p-1}$ divides $d_n$.  Since the $\beta_p$ are integers,
\eqref{comp:b}--\eqref{comp:c} are integers.  Their prime support is
contained in $\ell\leq2n+1$, without any assumption about the kernel's
rank.  For example,
\begin{align}
 12J_1&=36\gamma-48+30\log2+6\log3,
 \label{comp:n1}\\
 60J_2&=600\gamma-1125+760\log2+200\log3+20\log5.
 \label{comp:n2}
\end{align}

To prove positivity, set
$e_k=2H_k-2H_{n-k}-1/(n+1-k)$.  Then $e_0<0$, $e_n>0$, and
\[
 e_{k+1}-e_k=
 \frac2{k+1}+\frac1{n-k}+\frac1{n+1-k}>0
 \qquad(0\leq k<n).
\]
Since $\beta_{L+k}>0$ and $\alpha_M>0$, the simple residues have one
change from negative to positive, allowing a zero at the transition.
Their total is zero.  Every proper tail $T_j$ in \eqref{comp:c} is
therefore strictly positive.  It follows that $c_{n,\ell}\geq0$ and
\begin{equation}\label{comp:positive-log}
 \sum_\ell c_{n,\ell}\log\ell
   =d_n\sum_{j=L}^{M-1}T_j\log j>0.
\end{equation}

\subsection{Decay, arithmetic size, and the counting argument}

The elementary inequalities $-\log(xy)\geq1-xy\geq1-x$ imply
\[
 0<J_n\leq\int_0^1\!\int_0^1
 x^n(1-x)^{n-1}y^n(1-y)^n\,dx\,dy
 \leq 2\,16^{-n}.
\]
For the last bound, use
$[x(1-x)]^{n-1}\leq4^{-(n-1)}$ and
$[y(1-y)]^n\leq4^{-n}$, and integrate the remaining factor $x$.
The standard prime number theorem in the form
\[
 \log\operatorname{lcm}(1,\ldots,m)=m+o(m)
\]
gives $\log d_n=2n+o(n)$, proving \eqref{comp:smallness}.
Also $\beta_p\leq2\cdot4^n$,
$|\alpha_p|\leq2\cdot4^n(4H_n+1)$, and
$v_\ell((p-1)!)=O(n\log(n+1))$ uniformly in $p,\ell$.
Equations \eqref{comp:b}--\eqref{comp:c}, with at most $n+2$ terms,
prove \eqref{comp:height}.  Polynomial factors contribute only
$O(\log(n+1))$ to its logarithm.

We record the precise elementary counting step.  If an integer matrix
$D$ has $m<k$ rows and columns respectively and entries of absolute value
at most $B\geq1$, put
\begin{equation}\label{comp:box}
 H=\left\lfloor(3kB)^{m/(k-m)}\right\rfloor+1.
\end{equation}
The $(H+1)^k$ vectors in $\{0,\ldots,H\}^k$ have at most
$(2kBH+1)^m\leq(3kBH)^m$ different images under $D$.
Because $H^{k-m}>(3kB)^m$, two of those vectors have the same image.
Their difference is a nonzero integer vector $w$ with $Dw=0$ and
$\max_i|w_i|\leq H$.  This proof requires no rank hypothesis and no
nonvanishing hypothesis on any associated linear form.

Apply this to the matrix
\[
 D_N=(c_{n,\ell})_{
 \substack{\ell\leq4N+1\ {\rm prime}\\N\leq n\leq2N}},
\]
where coefficients outside their individual prime supports are zero;
primes index rows and $n$ indexes columns.  It has
$k=N+1$ columns and
$m=\pi(4N+1)=O(N/\log N)$ rows.  The preceding height estimate gives
$\log B=O(N)$, so \eqref{comp:box} gives
\[
 \log H\leq\log2+\frac{m}{k-m}\log(3kB)
          =O(N/\log N).
\]
Let $u$ be the resulting kernel vector.  Summing
\eqref{comp:integer-form} cancels the logarithms, with
$Q_N=\sum u_na_n$ and $P_N=\sum u_nb_n$.
Uniformly for $N\leq n\leq2N$,
$d_nJ_n\leq\exp(-cN+o(N))$; consequently
\[
 \left|\sum_{n=N}^{2N}u_nd_nJ_n\right|
 \leq(N+1)H\max_{N\leq n\leq2N}d_nJ_n
 \leq\exp(-cN+o(N)).
\]
This proves the existence and size statements in the theorem.

Finally, take the row vector with entries $\log\ell$.
By \eqref{comp:positive-log}, every coordinate of its product with $D_N$
is strictly positive.  Hence $D_Nu=0$ is impossible for a nonzero vector
whose coordinates all have the same weak sign.  The positivity of
$J_n$ therefore supplies no positivity of the compressed form.
To make the distinction still more explicit, append the two integer
rows $(a_n)$ and $(b_n)$ to $D_N$.  Their entries also have size
$\exp(O(N))$, and the row count is still $O(N/\log N)$.
The same counting argument produces a nonzero integer vector $v$ of
the same height with
\[
 D_Nv=0,\qquad\sum_{n=N}^{2N}v_na_n=0,
 \qquad\sum_{n=N}^{2N}v_nb_n=0.
\]
Thus $\sum v_nd_nJ_n=0$ exactly.  This completes the proof of
Theorem~\ref{comp:main} and identifies the missing step in any proposed
irrationality application: one must produce infinitely many compressed
forms that are provably nonzero.  No such claim is made here.

\section{Exact coordinates and finite-system obstructions for two reciprocal sums}
\label{qor:section}

The Erd\H{o}s--Borwein constant and the reciprocal Fibonacci sum are
\[
 E=\sum_{n\geq1}\frac{1}{2^n-1},\qquad
 \Psi=\sum_{n\geq1}\frac1{F_n},
 \qquad F_1=F_2=1,\quad F_{n+2}=F_{n+1}+F_n.
\]
We retain these complete sums.  They are values of the same elementary
meromorphic family.  The family has a first-order multiplicative
shift equation, but cannot be a component of a finite rational
power-Mahler system.  We prove that distinction and the associated
orbit identities explicitly.

This section consolidates statements T6, T27 and T28 of the supplied
HTML.  Their scalar and vector obstructions are one result and its
linear-algebra consequence, rather than independent discoveries.
The supplied corpus already contains the integer-base shift identity
in manuscript [030] \cite{Corpus030}, equation \texttt{p04:ci20:shift}; the jet formula
below follows by differentiating and iterating that type of identity.
The arguments are elementary specializations of established methods.
No historical priority or transcendence of either displayed constant
is asserted.

\subsection{One family, with the original constants preserved}

For a complex number $Q$ with $|Q|>1$, define
\begin{equation}\label{qor:L}
 L_Q(z)=\sum_{n\geq1}\frac{z^n}{Q^n-1},\qquad |z|<|Q|.
\end{equation}
Absolute convergence follows by comparison with
$\sum_{n\geq1}(|z|/|Q|)^n$.  Put $\phi=(1+\sqrt5)/2$.

\begin{lemma}[Exact value coordinates]\label{qor:coordinates}
The identities
\begin{equation}\label{qor:target-values}
 E=L_2(1),\qquad \Psi=\sqrt5\,L_{-\phi^2}(-\phi)
\end{equation}
hold with absolute convergence.  The functions satisfy the meromorphic
identity
\begin{equation}\label{qor:shift}
 L_Q(Qz)-L_Q(z)=\frac{z}{1-z}.
\end{equation}
Moreover,
\begin{equation}\label{qor:poles}
 L_Q(z)=\sum_{m\geq1}\frac{z}{Q^m-z}
\end{equation}
gives a meromorphic continuation to $\mathbb C$.  Its poles are exactly
$Q,Q^2,Q^3,\ldots$, each simple, with residue $-Q^m$ at $Q^m$.
\end{lemma}
\begin{proof}
The first equality in \eqref{qor:target-values} is the definition.
Binet's formula gives
\[
 \sqrt5 F_n=\phi^n-(-1)^n\phi^{-n},\qquad
 \frac{(-\phi)^n}{(-\phi^2)^n-1}
 =\frac{\phi^n}{\phi^{2n}-(-1)^n}
 =\frac1{\sqrt5 F_n}.
\]
Since $\phi<\phi^2$, both series representations converge absolutely.
For $|z|<1$, subtraction of the series proves
\eqref{qor:shift}.  Also, for $|z|<|Q|$,
\[
 \frac1{Q^n-1}=\sum_{m\geq1}Q^{-mn}.
\]
The absolute double sum is finite, so summing first in $n$ gives
\eqref{qor:poles}.  On any compact set, the tail of the latter sum
is bounded by a constant times $\sum_m|Q|^{-m}$.  It is therefore
normally convergent away from the listed poles.  Their absolute
values are distinct, so only one summand is singular at each pole;
its residue is $-Q^m$.  Analytic continuation now proves the shift
identity on its entire meromorphic domain.
\end{proof}

\subsection{Logarithmic derivatives and the exact orbit cocycle}

The product
\begin{equation}\label{qor:E-product}
 \mathcal E_Q(x)=\prod_{m\geq1}(1+x/Q^m)
\end{equation}
converges normally on compact sets and defines an entire function.
Its zeros are the simple points $-Q^m$, $m\geq1$.  Its logarithmic
derivative is
\begin{equation}\label{qor:R}
 R_Q(x)=\frac{\mathcal E'_Q(x)}{\mathcal E_Q(x)}
       =\sum_{m\geq1}\frac1{Q^m+x}.
\end{equation}
These assertions follow from locally uniform convergence of
$\sum_m|x/Q^m|$ and of the differentiated logarithms away from
their zeros.  In particular,
\begin{equation}\label{qor:R-values}
 L_Q(z)=zR_Q(-z),\qquad
 E=R_2(-1),\qquad
 \Psi=-\sqrt5\phi R_{-\phi^2}(\phi).
\end{equation}

For an integer $s\geq0$ write $R_Q^{(s)}$ for the ordinary $s$th
derivative.  Set $c_s=(-1)^ss!$, and for every integer $k$ define
the rational function
\begin{equation}\label{qor:C}
 C_{k,s}(x)=
 \begin{cases}
 \displaystyle\sum_{j=0}^{k-1}
       \frac{Q^{j(s+1)}}{(1+Q^jx)^{s+1}},&k>0,\\[6pt]
 0,&k=0,\\[3pt]
 \displaystyle-\sum_{j=k}^{-1}
       \frac{Q^{j(s+1)}}{(1+Q^jx)^{s+1}},&k<0.
 \end{cases}
\end{equation}

\begin{theorem}[All derivative orders and both orbit directions]
\label{qor:jets}
For $|Q|>1$, $s\geq0$ and $k\in\mathbb Z$, the meromorphic identity
\begin{equation}\label{qor:jet-cocycle}
 Q^{k(s+1)}R_Q^{(s)}(Q^kx)
       =R_Q^{(s)}(x)+c_s C_{k,s}(x)
\end{equation}
holds.  Suppose additionally that $Q,x\in\overline{\mathbb Q}$
and that the values under consideration are finite.  The regular
values of a fixed derivative order on this orbit belong to
\[
 \operatorname{span}_{\overline{\mathbb Q}}
       \{1,R_Q^{(s)}(x)\}.
\]
Consequently all regular orbit values of orders $0\leq s\leq S$
span a space of dimension at most $S+2$ over
$\overline{\mathbb Q}$.
\end{theorem}
\begin{proof}
Separating the first factor of the product gives
\[
 \mathcal E_Q(Qx)=(1+x)\mathcal E_Q(x),\qquad
 QR_Q(Qx)=R_Q(x)+\frac1{1+x}.
\]
Differentiating $s$ times yields
\[
 Q^{s+1}R_Q^{(s)}(Qx)
       =R_Q^{(s)}(x)+\frac{c_s}{(1+x)^{s+1}}.
\]
For $j\in\mathbb Z$, substitute $Q^jx$ and multiply by
$Q^{j(s+1)}$.  The left side minus its first right-side term is
the difference of consecutive terms of the sequence
$Q^{j(s+1)}R_Q^{(s)}(Q^jx)$.  Telescoping forward or backward proves
\eqref{qor:jet-cocycle}.

For finite evaluation at algebraic $Q,x$, each term of
$C_{k,s}(x)$ is algebraic provided its denominator is nonzero.
This is automatic when both endpoint values in
\eqref{qor:jet-cocycle} are regular.  Indeed a vanishing denominator
means $x=-Q^{-j}$ for an intermediate index $j$; in either direction
one endpoint then has the form $-Q^m$ with $m\geq1$ and is a pole.
Thus every regular value is an algebraic multiple of the base
value plus an algebraic number.  Taking the indicated spans proves
the two bounds.
\end{proof}

\begin{remark}[Affine information and the pole qualification]
Modulo the additive subspace $\overline{\mathbb Q}$, the normalized
values $Q^{k(s+1)}R_Q^{(s)}(Q^kx)$ all have the same class.  This is
the precise sense in which a fixed derivative order supplies at most
one new orbit coordinate.  It does not assert that the coordinate
is transcendental, or that all determinants assembled in arbitrary
ways from orbit data vanish.

The qualification about poles matters for the original targets.
For $E=R_2(-1)$ every forward point $-2^k$, $k\geq1$, is a pole.
Its regular backward orbit is $-2^{-k}$, $k\geq0$; for example,
\[
 R_2(-1/2)=2R_2(-1)-2=2E-2.
\]
For $Q=-\phi^2$ and $x=\phi$, every integer orbit point is regular:
an equality $Q^k\phi=-Q^m$ would imply
$\phi^{2k+1}=\phi^{2m}$, impossible for integers $k,m$.
\end{remark}

\subsection{A complete obstruction to finite power-Mahler embedding}

A power-Mahler equation uses the substitution $z\mapsto z^d$ for
an integer $d\geq2$.  It is distinct from the multiplicative shift
$z\mapsto Qz$ in \eqref{qor:shift}.  The pole sets make that
distinction exact here.

\begin{theorem}[Independence of all power iterates]
\label{qor:no-mahler}
Let $Q\in\mathbb C$ satisfy $|Q|>1$, and let $d\geq2$ be an integer.
For every $r\geq0$, the functions
\[
 1,\ L_Q(z),\ L_Q(z^d),\ldots,L_Q(z^{d^r})
\]
are linearly independent over $\mathbb C(z)$.  Equivalently, if
$a_0,\ldots,a_r,b\in\mathbb C(z)$ and
\begin{equation}\label{qor:relation}
 \sum_{j=0}^r a_j(z)L_Q(z^{d^j})=b(z)
\end{equation}
holds as an analytic or meromorphic germ, then every coefficient
and $b$ vanish identically.
\end{theorem}
\begin{proof}
Assume some $a_j$ is nonzero, and let $r$ now denote its greatest
index.  For every positive integer $m$ with $d\nmid m$, choose any
$z_m$ satisfying
\[
 z_m^{d^r}=Q^m.
\]
The numbers $z_m$ are nonzero and have distinct absolute values
$|Q|^{m/d^r}$.  By Lemma~\ref{qor:coordinates}, the highest iterate
$L_Q(z^{d^r})$ has a simple pole at each $z_m$; the derivative of
$z^{d^r}$ there is nonzero.

No lower iterate has a pole there.  Otherwise for some $j<r$
and $n\geq1$ one would have $z_m^{d^j}=Q^n$.  Raising to
$d^{r-j}$ gives
\[
 Q^m=Q^{nd^{r-j}}.
\]
Taking absolute values, using $|Q|>1$, implies
$m=nd^{r-j}$ and hence $d\mid m$, a contradiction.
This argument also covers $r=0$, for which there are no lower
iterates.

Exclude the finitely many zeros and poles of the nonzero rational
function $a_r$ and the finitely many poles of $b,a_0,\ldots,a_{r-1}$.
Infinitely many $z_m$ remain.  At any one of them the left side of
\eqref{qor:relation} has a pole from exactly its highest term,
whereas the right side is regular.  The meromorphic continuation
of the germ identity is therefore impossible.  All $a_j$ vanish,
and then so does $b$.
\end{proof}

\begin{remark}[Negative and complex bases are included]
The proof uses absolute values to compare powers of $Q$ and never
requires $Q>0$.  In particular it applies to $Q=-\phi^2$ without
a choice of logarithm.  An alternative explicit construction for
real $Q$ and $r\geq1$ is to choose a primitive $d^r$th root
$\zeta$ and set $z_t=\zeta|Q|^{2t}$.  The highest iterate has its
pole at $Q^{2td^r}$.  A pole in a lower iterate would require the
even index $n=2td^j$, hence a positive pole $Q^n$, while
$\zeta^{d^j}\ne1$.  This gives the same contradiction.
\end{remark}

\begin{corollary}[No finite rational affine system]
\label{qor:no-vector}
There is no finite vector $Y$ of analytic or meromorphic germs,
having $L_Q$ as a component, that satisfies
\[
 Y(z^d)=A(z)Y(z)+B(z),
 \qquad A\in\operatorname{Mat}_n(\mathbb C(z)),\quad
 B\in\mathbb C(z)^n.
\]
The assertion holds even without assuming that $A$ is invertible.
\end{corollary}
\begin{proof}
Adjoin the constant component $1$.  Iteration of the displayed
system expresses every component of $Y(z^{d^j})$ in the fixed
$\mathbb C(z)$-vector space generated by $1,Y_1(z),\ldots,Y_n(z)$.
It follows that all the iterates of the component $L_Q$ lie in a
space of dimension at most $n+1$.  Sufficiently many of them are
linearly dependent, contrary to Theorem~\ref{qor:no-mahler}.
\end{proof}

\subsection{What these statements leave unresolved}

The value identities \eqref{qor:target-values} and
\eqref{qor:R-values} preserve the two named constants exactly.
The orbit calculation controls information at points related by
one fixed multiplicative shift.  The last theorem excludes a
finite rational power-Mahler system containing the specified
function.  It does not exclude a different function that happens
to take the same value at one point, and it does not exclude other
arithmetic methods.

No value-transcendence inference follows merely from functional
independence.  For example, the same family has the algebraic
value $L_Q(0)=0$, and for algebraic $Q$ its logarithmic derivative
has
\[
 R_Q(0)=\sum_{m\geq1}Q^{-m}=\frac1{Q-1}
       \in\overline{\mathbb Q}.
\]
Proving that $R_2(-1)$ or $R_{-\phi^2}(\phi)$ is transcendental
requires a theorem for those particular values that is not supplied
by the pole argument or the jet cocycle.  The unconditional output
of this section is a precise restriction on two proposed functional
routes, with full proofs and unchanged target values.

\section{Consequences and remaining research tasks}\label{sec:limits}
The conclusions have different logical strengths. For the quartic
Salem number and the specified prime set, the condition $260\mid h$
is necessary and sufficient, and the density table is complete.
The repeated-block and pole arguments are also unconditional theorems
with the exact domains stated in their hypotheses.

For Euler's constant, cross-index elimination removes every logarithm
with controlled integer coefficients. However, appending only two more
rows to the same integer matrix already produces comparably short
vectors whose resulting form is exactly zero. An irrationality
argument therefore needs a mechanism that proves nonvanishing for
infinitely many selected forms; vector nontriviality and positivity of
the individual integrals cannot provide it.

For the two reciprocal sums, the meromorphic family and the original
values are unchanged. Its multiplicative-shift jets contain at most
one new coordinate per derivative order, and its power iterates cannot
lie in any finite rational affine system. These conclusions restrict
specified approaches and supply no assertion about every possible
representation of the same value.

Finally, the twin-prime partial-product theorem concerns the exact
unweighted rational truncations in its statement. Weighted products,
other sequences of approximants, and unrelated methods are outside
that theorem. Keeping each conclusion at its proved scope is necessary
for this integrated manuscript to be reusable as a research starting
point.
\clearpage\appendix
\section{A self-contained exact arithmetic certificate}
\label{app:certificate}

The listing below is part of this manuscript.  It can be copied unchanged
to a text file and run with Python~3; it needs only the standard library.
It checks the trace period, the unsafe-divisor witnesses, every exponent
class modulo the trace period, the complete rational density table, and
an independent quotient-ring computation of the matrix order and every
density for all $780$ exponent classes.  Thus no
separate code file is required to reproduce the finite part of the proof.
The strict correction window and the equidistribution argument in the
main text supply the infinite-index implications; those analytic results
are not inferred from this finite calculation.

The quotient-ring computation uses
\[
 X^4=X^3+X^2+X-1
 \quad\text{in }(\mathbb Z/30\mathbb Z)[X]/(P(X)).
\]
Multiplication by $X$ in the basis $1,X,X^2,X^3$ is exactly the companion
matrix displayed in the main text.  Consequently its order is the order
computed below.  For a representative $r_0+r_1X+r_2X^2+r_3X^3$,
the trace of its multiplication matrix is
$4r_0+r_1+3r_2+7r_3$ modulo $30$; this follows by taking the traces of
$I,C,C^2,C^3$.  The recurrence state transition is invertible, so the
first repeated initial four-term state proves the exact trace period for
all indices.

\begin{lstlisting}[language=Python,basicstyle=\ttfamily\footnotesize,
  columns=fullflexible,keepspaces=true,breaklines=true,showstringspaces=false,
  frame=single,numbers=left,numberstyle=\tiny,
  caption={Exact certificate for the quartic classification.},
  label={lst:salem-certificate}]
from fractions import Fraction as F
from math import gcd

M = 30
T = [4, 1, 3, 7]
for n in range(4, 264):
    T.append((T[-1] + T[-2] + T[-3] - T[-4]) % M)
period = next(k for k in range(1, 261)
              if T[k:k+4] == T[:4])
assert period == 260

Q = (-1, 0, 1, 2)
safe = {s for s in range(M)
        if all(gcd(s-q, M) > 1 for q in Q)}
assert safe == {4, 5, 26, 27}
divisors = [g for g in range(1, 261) if 260 % g == 0]
expected = {
    1: F(541,1560), 2: F(241,780), 4: F(14,39),
    5: F(97,312), 10: F(15,52), 13: F(49,120),
    20: F(8,39), 26: F(7,20), 52: F(2,5),
    65: F(5,24), 130: F(1,12), 260: F(0)
}
assert divisors == list(expected)
assert [T[g] for g in divisors[:-1]] == [
    1, 3, 7, 16, 18, 1, 22, 9, 19, 16, 24
]
weights = {-1: F(1,3), 0: F(1,6),
           1: F(1,6), 2: F(1,3)}
for g in divisors:
    orbit = range(0, period, g)
    value = sum((weights[q] for r in orbit for q in Q
                 if gcd(T[r]-q, M) == 1), F(0))
    value /= len(orbit)
    assert value == expected[g]
    assert (all(T[r] in safe for r in orbit)) == (g == 260)
    print(g, value)

for h in range(1, 261):
    residues = {h*n % period for n in range(period)}
    covered = all(T[r] in safe for r in residues)
    assert covered == (h % 260 == 0)
assert min(v for v in expected.values() if v) == F(1,12)
assert [g for g, v in expected.items() if v == F(1,12)] == [130]

# Independent multiplication in (Z/30Z)[X]/(P).
def mul(a, b):
    c = [0] * 7
    for i in range(4):
        for j in range(4):
            c[i+j] += a[i] * b[j]
    for k in range(6, 3, -1):
        v = c[k]
        c[k-1] += v
        c[k-2] += v
        c[k-3] += v
        c[k-4] -= v
    return tuple(v % M for v in c[:4])

one, x = (1,0,0,0), (0,1,0,0)
r = one
orbit = [one]
for n in range(1, 781):
    r = mul(r, x)
    if n == 260:
        assert r == (21,20,20,20)
    if n == 520:
        assert r == (11,10,10,10)
    assert (r == one) == (n == 780)
    if n < 780:
        orbit.append(r)
assert len(set(orbit)) == 780

# Matrix traces from the quotient orbit, not from T.
basis_traces = (4, 1, 3, 7)
matrix_traces = [sum(a*b for a, b in zip(r, basis_traces))
                 % M for r in orbit]
for h in range(1, 781):
    length = 780 // gcd(h, 780)
    traces = [matrix_traces[h*k % 780]
              for k in range(length)]
    value = sum((weights[q] for s in traces for q in Q
                 if gcd(s-q, M) == 1), F(0)) / length
    assert value == expected[gcd(h, 260)]
    assert all(s in safe for s in traces) == (h % 260 == 0)
print("PASS: trace period 260, matrix order 780,")
print("all 780 matrix-orbit density classes verified.")
\end{lstlisting}

All numerical quantities in the listing are integers or reduced rational
numbers.  No numerical approximation to $\tau$, its conjugates, or their
angles is used.  In particular, the listing does not test primality and
does not identify the terms coprime to $30$ with prime terms.

\clearpage

\section{Source concordance for the selected results}\label{app:sources}
The labels T1--T60 refer to the supplied HTML record, not to new
theorem numbers in this paper. Corpus identifiers refer to the numbered
files at the immutable snapshot in Section~\ref{sec:introduction}.
Additional below means that the same substantive statement was not
located in that snapshot; it is not a worldwide priority assertion.

\small
\begin{longtable}{@{}>{\raggedright\arraybackslash}p{.12\textwidth}>{\raggedright\arraybackslash}p{.38\textwidth}>{\raggedright\arraybackslash}p{.44\textwidth}@{}}
\toprule
Source & Retained mathematical content & Relation to the corpus\\
\midrule
\endfirsthead
\toprule Source & Retained mathematical content & Relation to the corpus\\
\midrule
\endhead
T1--T3 &
No macroscopic repeated blocks in the prime-count residue word,
prime-gap word and divisor-count word. &
Additional block statements. Manuscript [003] contains the prime
codings, while [030] studies related Lambert functions. We state
block theorems directly, avoiding an unverified transfer to the
binary digits after carrying.\\[5pt]
T9 &
Simple-pole logarithmic positivity and the two-location
classification. &
Extends the complete pole transform in [001],
\texttt{prop:complete-pole-transducer}. The transform itself is
reproduced and credited.\\[5pt]
T13 &
Reduced denominator versus tail for the raw twin-prime
Euler product. &
No same substantive result located. The proof uses the elementary
two-adic denominator bound and the prime number theorem.\\[5pt]
T55 &
Short cross-index logarithm elimination for the asymmetric
Euler integral, including the zero-form obstruction. &
Synthesis: the exact integral and its arithmetic occur in [001],
\texttt{thm:Jn-asymptotic-integerization}; the moving-window counting
mechanism already occurs in [021] for a different integral family.\\[5pt]
T6, T27, T28 &
Fixed-$Q$ Lambert poles, the all-jet orbit identity, and no finite
rational power-Mahler embedding. &
The integer-base shift equation already occurs in [030],
\texttt{p04:ci20:shift}. The fixed-$Q$ pole proof differs from its
variable-base Lambert theorem. Scalar and vector obstructions are
combined into one result and its corollary.\\[5pt]
T60 and the preceding Salem note &
The strict correction mechanism, prescribed-exponent criterion,
quartic equivalence $260\mid h$, and all twelve exact densities. &
The norm window, independent angles, existence construction and sharp
prime-cover sets occur in [005]/[031]. The preceding note adds the
fixed-exponent classification and its density table; these are
retained in full.\\
\bottomrule
\end{longtable}
\normalsize

The comparison was performed across all 49 TeX files, rather than by
matching titles alone. Nineteen of T41--T60 have substantive matches in
the corpus; T55 is a literal-family synthesis of existing methods.
Other omitted items include already incorporated chamber-moment and
denominator results in [011], and the conditional Mills identities in
[010]/[031]. Their omission is an editorial decision to avoid
reprinting existing material as an additional result.

Two scope repairs are reflected in the mathematical text. First,
the forward orbit of $R_2(-1)$ consists of poles, so finite-value orbit
identities for the Erd\H{o}s--Borwein constant use its backward orbit.
Second, the quartic trace period is $260$, whereas the companion
matrix order modulo $30$ is $780$. The embedded certificate checks
both quantities separately and confirms the same densities by both
representations.

\clearpage
\begin{thebibliography}{99}
\bibitem{HTMLRecord}
Supplied research record,
\emph{Transcendence full-proof research audit} (HTML, Korean),
statements T1--T60, supplied uploaded revision.
Primary working material supplied for this
integration; the labels in Appendix~\ref{app:sources} identify the
retained statements.
\bibitem{SalemNote}
\emph{Fixed-Exponent Prime Covers for Salem Power Floors:
An exact quartic exponent and density classification},
research note prepared for KEIO, 24 October 2026.
The complete mathematical argument is retained in this manuscript.
\bibitem{Corpus}
QEEDDD, \emph{Verifiable Mathematical Reasoning Corpus},
49 numbered TeX research manuscripts, immutable revision
\texttt{f70e5fedabf7b90313f703eb85f7dff72a055541},
checked 24 October 2026 (Korea Standard Time).
\href{https://github.com/QEEDDD/Verifiable-Mathematical-Reasoning-Corpus/tree/f70e5fedabf7b90313f703eb85f7dff72a055541}{Repository snapshot}.
The following numbered items refer to this same revision.
\bibitem{Corpus001}
QEEDDD, corpus [001],
\emph{Exponential-Integral Periods across Arithmetic, Pade, Hankel,
Reciprocal Geometry, and Frobenius Framing}, especially the complete
pole transform and the asymmetric Euler-integral family.
\bibitem{Corpus003}
QEEDDD, corpus [003],
\emph{Integrated Structures and Boundary Spectra in Sparse Prime
Codings}.
\bibitem{Corpus005}
QEEDDD, corpus [005],
\emph{Composite Integer Parts of Powers of Algebraic Units
(Pure Cubic Fields, Salem Numbers, and a K3 Application)}.
\bibitem{Corpus021}
QEEDDD, corpus [021], \emph{Arithmetic Alternatives for Euler’s Constant}.
Moving-window integer-kernel elimination of logarithmic coefficients.
\bibitem{Corpus030}
QEEDDD, corpus [030],
\emph{Arithmetic and Observation: Classifications, Recovery, and
Transcendence Barriers. Unified Edition of Twelve Successor Manuscripts},
especially equation \texttt{p04:ci20:shift}.
\bibitem{Corpus031}
QEEDDD, corpus [031],
\emph{Prime-Representing Constants and Related Structures
(Encodings, Arithmetic, Geometry, and Modular Obstructions)},
section on sharp finite prime covers for algebraic power floors.
\bibitem{HIKK}
S.~Hahn, D.~Ismailescu, G.~Kim and M.~Kim,
\emph{Explicit Algebraic Numbers all whose Integer Parts of Powers
are Composite}, arXiv:2608.05309v1 (2026), especially Section~7.
\url{https://arxiv.org/abs/2608.05309}.
\bibitem{ArithDavenport}
H.~Davenport, \emph{Multiplicative Number Theory}, third edition,
revised by H.~L. Montgomery, Graduate Texts in Mathematics 74,
Springer, 2000.
\bibitem{ArithHW}
G.~H. Hardy and E.~M. Wright, \emph{An Introduction to the Theory
of Numbers}, sixth edition, revised by D.~R. Heath-Brown and
J.~H. Silverman, Oxford University Press, 2008.
\end{thebibliography}
\end{document}
